\chapter{Parameterstudien, Berichte und Surrogate}\label{sec:app-studies}
\index{Study}\index{Wiederaufnahme}\index{Surrogat}
\section{Eine Studie bewusst unterbrechen und fortsetzen}
Studien-Schema 1 besitzt \path{study_schema_version}, \path{model} und
\path{parameters}. Jede Parametertabelle enthält Name, vorhandenen Zielpfad,
Einheit und eindeutige endliche Werte. Es wird das kartesische Produkt
in deterministischer Reihenfolge ausgeführt; keine Zufallsziehung oder
Zip-Semantik. Vektoren/Tensoren werden als Ganzes mit unveränderter Form
ersetzt. Neue Felder/Topologien erzeugt ein Parameterpfad nicht.
Nicht-SI-Ziele vor dem Studium explizit nach SI umrechnen; Solver-/Pfadwerte
verwenden die dimensionslose Einheit \path{1}. Studien sind lokal und
sequenziell, ausschließlich die unterstützten stationären Model-Physiken.

\begin{lstlisting}
python -I examples/parameter_study/run.py --output results/study --max-cases 1
python -I examples/parameter_study/run.py --output results/study
python -I examples/parameter_study/run.py --output results/study
python -I examples/parameter_study/report.py --study results/study --output results/table
\end{lstlisting}
Der erste Aufruf begrenzt bewusst auf einen neuen Versuch. Der zweite beendet
die zwei Fälle; der dritte prüft Wiederverwendung. Ein eigenes Resume-Flag
gibt es nicht. Im \path{catalog.json} sollen nach Abschluss beide Fälle
\path{succeeded} mit \path{attempts=1} stehen; der dritte Aufruf erhöht
diese Zahl nicht. Erwartete Maximaltemperaturen sind 301 K und 300.5 K.
Ein Kontrolllauf ohne \path{--max-cases} in einem frischen Ziel liefert
unabhängig dieselben Größen.

Fall-ID ist SHA-256 der kanonischen benannten Parameterwerte und Einheiten.
Fingerprints binden aufgelöstes Modell, tatsächliche Netzbytes, Solveroptionen
und Versionsvertrag. Ein Erfolg wird nur bei übereinstimmenden Fingerprints
und gültigen Bündelhashes übernommen. Andere NumPy-/SciPy-/Toolbox-Versionen
können Neuberechnung auslösen. Unveröffentlichte Solveränderungen benötigen
einen neuen Vertrag oder ein neues Ausgabeziel. Änderungen am Basismodell
invalidieren Hashes bei gleichbleibender Parameterfall-ID. Eingaben während
eines Laufs unverändert lassen; ein neues Studienmanifest braucht ein neues
\path{Study}-Objekt. Dateipfade bleiben innerhalb der erlaubten Modellbäume.

Fehlfälle sind sichtbar mit Fehler und Versuchszahl; sie stoppen nicht
sämtliche anderen Fälle. Nach einem Prozessabbruch erneut denselben Aufruf
starten und den Katalog prüfen. Ein liegengebliebener laufender Versuch kann
neu versucht werden; erfolgreiche geprüfte Bündel bleiben wiederverwendbar.
Keine manuelle Katalogstatusänderung als Ersatz für gültige Artefakte.
Die Wärme- und Piezorezepte \ref{recipe:study-recipes} prüfen zusätzlich
Unterbrechung, zweiten unabhängigen Direktlauf, Fingerprintänderung und
absichtlich unzulässige Materialien. Der Piezoaufruf lautet:
\begin{lstlisting}
python -I examples/study_recipes/run.py piezo --output results/piezo-study
\end{lstlisting}
Das jeweilige \path{verification.json} enthält Referenzen, Istwerte,
Toleranzen, Runtime und tatsächlichen Paketimportpfad.

\section{Berichte aus bestehenden Bündeln}
\path{Quantity} definiert Namen, Rohfeld-/Metrikquelle, Einheit und
gegebenenfalls Reduktion. Point-Reduktionen berücksichtigen aktive Knoten,
nicht Exportfüllwerte. Für Bilanzprüfungen ist die absolute Toleranz explizit;
kleine Werte werden nicht aus optischen Gründen auf null gesetzt.
\path{export_study_table} liest den Katalog und die gesicherten Bündel,
rechnet keine neue FEM-Lösung und erzeugt ein neues Zielverzeichnis mit
\path{values.csv}, \path{metadata.json}, \path{report.md}.
\path{read_study_table} prüft Schema, CSV-Hash, Einheiten und Metadaten.
\path{table_schema_version=1} ist unabhängig vom Studien- und Bündelschema.

Die Tabellen bewahren Fall-ID, Parameter, Status, fehlende Werte, Quellen-
und Ergebnisfingerprints. Fehlende oder fehlgeschlagene Fälle werden nicht
zu erfolgreichen Nullwerten. Der einfache Bericht enthält Maximaltemperatur
in K und vorzeichenbehaftete Leistungsbilanz in W. Der Surrogatbericht
zusätzlich Quellenleistung und Einheitenvertrag. Für einen eigenen Messpunkt
oder ein Integralmittel zunächst eine passende, unterstützte \path{Quantity}
festlegen; keine ungeprüfte Mittelung von Schwerpunkt- und Knotenwerten.

\section{Surrogat trainieren, speichern, wieder laden}
Rezept \ref{recipe:surrogate} erstellt acht FEM-Fälle eines Einheits-Tet4.
Sechs Quellenwerte trainieren ein skaliertes Polynom Grad 1 mit
\path{ridge=0}; die getrennten Testwerte 3 und 7 W/m$^3$ bleiben zurückgehalten.
Die analytische Referenz lautet $T_{\max}=300+Q/8$ K und $P_Q=Q/6$ W.
Das Skript prüft die FEM-Zielwerte und Vorhersagen an beiden Testfällen mit
$10^{-10}$ K bzw. $10^{-12}$ W und lädt das gespeicherte JSON erneut.
\begin{lstlisting}
python -I examples/surrogate_heat/run.py --output results/surrogate
python -I -c "from fem_toolbox.surrogates import PolynomialSurrogate; m=PolynomialSurrogate.load('results/surrogate/surrogate.json'); print(m.predict([5.0]).values.tolist())"
\end{lstlisting}
Erwartet werden \path{[[300.625, 0.8333333333333333]]} bis zu den angegebenen
Rundungstoleranzen. Die Reihenfolge der Ziele folgt dem Datensatz, hier K
und W; die Rohzahl trägt nicht automatisch eine beliebige physikalische Einheit.
\path{SurrogateData.from_table} wählt explizite Parameter-/Zielspalten.
\path{DataSplit} bindet Trainings-/Test-Fall-IDs ohne Überlappung;
fehlgeschlagene Zeilen und ungültige Eingaben sind kein Trainingsmaterial.
Der Bericht enthält Datenherkunft, Split, Skalierung und getrennte Fehler.
\path{surrogate_schema_version=1} und Dateiintegrität werden beim Laden geprüft.

Die Parameterbox ist $Q\in[0,10]$ W/m$^3$. Vorhersage außerhalb wird
standardmäßig abgelehnt; \path{allow_extrapolation=True} muss bewusst gesetzt
werden und kennzeichnet extrapolierte Zeilen. Ein kleiner Testfehler im
Inneren liefert keine Extrapolationsgarantie. Konstant gewählte Eingänge,
unzureichender Rang und inkonsistente Einheiten dürfen nicht durch geschönte
Berichte verdeckt werden. Konstante Zielgrößen sind möglich; deren relative
Fehlerinterpretation benötigt den dokumentierten Maßstab.

\section{Noch offener Unsicherheitsworkflow}
M9-Issue 64 ist nicht implementiert. Es gibt hier keine ausführbare
Unsicherheits-API, keine angenommenen Messunsicherheiten und keine aus
fehlenden Messungen abgeleiteten Konfidenzintervalle. Der aktuelle
Surrogat-Testfehler ist numerischer Approximationsfehler, keine experimentelle
Unsicherheit. Nach \#64 sind Benutzerrezept mit ausdrücklich synthetischen
Lerndaten, API-Referenz und Quellen-/Rezeptprüfungen gemeinsam nachzuziehen.
Unabhängige Piezo-Messdaten samt Nutzungsrechten und Metadaten bleiben die
Voraussetzung der experimentellen M9-Validierung (\#63).
